% ============================================================
% DADOS: TERMO DE CONSENTIMENTO PARA ACESSO A PRONTUÁRIO
% Arquivo: dados/dados-termo-consentimento-prontuario.tex
%
% Peça auxiliar, fora dos anexos C-Y.
% Preencha conforme o caso concreto.
% ============================================================

% ------------------------------------------------------------
% Identificação do militar que autoriza.
% Por padrão, usa os dados gerais do sindicado.
% ------------------------------------------------------------
\newcommand{\tcpNome}{\getSindicadoNome}
\newcommand{\tcpPostoGraduacao}{\getSindicadoPosto}
\newcommand{\tcpIdentidade}{\getSindicadoIdentidade}
\newcommand{\tcpCpf}{\getSindicadoCPF}

% Exemplo: "Militar", "Soldado do Efetivo Variável", "Sd EV", etc.
\newcommand{\tcpQualificacao}{Militar}

% Exemplo: "sindicado", "interessado" etc.
\newcommand{\tcpCondicaoProcessual}{sindicado}

% ------------------------------------------------------------
% Objeto médico delimitado.
% Evite autorização ampla demais. Descreva somente o necessário.
% ------------------------------------------------------------
\newcommand{\tcpObjetoMedico}{queixas, atendimentos, exames e registros médicos relativos a dores ou lesões nos joelhos, no contexto dos fatos apurados}

% ------------------------------------------------------------
% DATA DO TERMO
%
% Opção 1, padrão: usa a data geral de dados-sindicancia.tex
% \newcommand{\tcpLocalData}{\localData}
%
% Opção 2, recomendada quando o termo tiver data própria:
% edite diretamente a linha abaixo.
% ------------------------------------------------------------
\newcommand{\tcpLocalData}{Petrolina/PE, 01 de junho de 2026}

% ------------------------------------------------------------
% Linha complementar abaixo da assinatura.
% Exemplo: "IDT nº 000000000 ÓRGÃO/UF" ou "CPF nº ..."
% ------------------------------------------------------------
\newcommand{\tcpIdentificacaoAssinatura}{IDT nº \tcpIdentidade}
